\section{Introduction}\label{sec:introduction}
% CHECKLIST-BEGIN: introduction
% P:introduction-carbon-capture
Reducing CO$_2$ emissions from fossil-fuel combustion is an important part of climate-change mitigation \cite{Harun2012a}.
Post-combustion capture separates CO$_2$ from flue gas before atmospheric release, providing a route to emissions reduction at power-generation and industrial sources \cite{Freguia2003,Harun2012a}.
Aqueous monoethanolamine (MEA) absorption is an extensively studied reference process for this separation, combining reactive solvent chemistry with packed-column contacting \cite{Freguia2003,Mores2012}.
Its established experimental and modeling basis makes it useful for evaluating how thermodynamic and transport descriptions affect capture predictions.

% P:introduction-process-need
The practical value of solvent-based capture depends on its separation performance, energy requirements, and ability to accommodate changes in gas and solvent conditions \cite{Harun2012a,Rodriguez2014}.
Solvent regeneration contributes to the process energy demand, while absorber operation determines the amount of CO$_2$ transferred to the rich solvent \cite{Freguia2003}.
Predicting capture together with the axial temperature profile is therefore necessary to connect solvent circulation and contacting conditions to the coupled mass- and heat-transfer response.
An absorber model supports these decisions when the effects of its physical assumptions can be distinguished from numerical error.

% P:introduction-01
MEA absorption couples chemical reaction, phase equilibrium, interfacial transport, and heat release in a countercurrent packed column.
Rate-based models resolve these interactions and support the interpretation of pilot-scale capture and temperature profiles \cite{Freguia2003,Zhang2009,Mores2012,Razi2013,Chinen2018}.
The calculated performance depends on the thermodynamic description, the liquid-film reaction model, and the numerical treatment of the mixed inlet conditions.
Controlled comparisons identify which modeling choices control the predicted response.

% P:introduction-02
Thermodynamics determines the distribution of dissolved CO$_2$ among molecular and ionic species, the interfacial driving force, and the energy released during absorption.
Changes in these quantities can affect both total capture and the axial temperature profile.
A column comparison therefore connects the liquid equilibrium description to its consequences for transfer and heat release.

% P:introduction-03
The liquid-film model determines how the thermodynamic driving force produces a transfer rate.
Its treatment of reaction and diffusion can redistribute absorption even when the thermodynamic model and feed conditions are unchanged.
Separating these effects requires a comparison at the same thermodynamic model, with each transfer model determining its own axial state.

% P:introduction-04
Countercurrent operation places gas and liquid inlet conditions at opposite ends of the packed bed.
The resulting boundary-value problem couples material and energy balances through temperature-dependent equilibrium and transport properties.
Numerical accuracy provides the basis for interpreting differences between physical models, while computational cost determines the effort required to make those comparisons.

\subsection{Literature review}\label{subsec:literature-review}
\subsubsection{Rate-based absorption and experimental evaluation}
% P:introduction-05
The rate-based literature provides the physical basis for coupling reaction, mass transfer, and heat release in aqueous MEA absorbers.
Early column studies established the importance of simultaneous material and energy balances, and subsequent work extended the models to solvent composition, gas loading, packing behavior, and pilot-scale operation \cite{Freguia2003,Zhang2009,Mores2012}.
Process studies further connect absorber behavior to load changes, control, and capture-process design \cite{Moioli2012,Harun2012a,Greer2008,Rodriguez2014}.
These applications require both integral predictions, such as capture, and local predictions, such as the position and magnitude of the temperature bulge.

% P:introduction-06
Pilot-plant studies connect the individual property and transport descriptions to measured absorber performance \cite{Razi2013,Chinen2018,Morgan2018a,Morgan2020}.
The National Carbon Capture Center (NCCC) studies combine operating-condition records with capture and temperature measurements, including the MEA-4 experimental-design campaign \cite{Morgan2018}.
Capture integrates transfer over the bed, whereas the axial temperature profile constrains the distribution of absorption, water evaporation, and interphase heat exchange.
Evaluation across operating cases tests these responses over changes in solvent circulation, gas throughput, and lean loading.
This provides a stronger basis for interpreting a model comparison than agreement at one operating point.

\subsubsection{Thermodynamic descriptions and their column implications}
% P:introduction-07
Electrolyte nonrandom two-liquid (eNRTL) models describe liquid nonideality through excess Gibbs energy and a specified species and reaction system.
For aqueous MEA, Zhang et al. combined vapor--liquid equilibrium, caloric, and speciation information to determine interaction parameters and ionic standard-state properties \cite{Zhang2011}.
That treatment connects the equilibrium driving force to both species distribution and thermal properties.
Its comparison with an electrolyte equation of state therefore concerns the complete thermodynamic description and its parameterization, as well as the choice of free-energy model.

% P:introduction-thermodynamic-literature
Statistical associating fluid theory (SAFT) provides a molecular description of association alongside the reaction-equilibrium calculation.
Baygi and Pahlavanzadeh examined perturbed-chain SAFT (PC-SAFT) association schemes for aqueous MEA \cite{Baygi2015}.
Najafloo and Zarei studied CO$_2$ solubility in aqueous MEA with Huang--Radosz SAFT (SAFT-HR) \cite{Najafloo2018}.
Electrolyte perturbed-chain statistical associating fluid theory (ePC-SAFT) adds ionic contributions to the associating-fluid description \cite{Gross2001,Cameretti2005}.
Cleeton et al. applied ePC-SAFT to competitive acid-gas absorption in aqueous methyldiethanolamine (MDEA), and B\"ulow et al. considered sour-gas absorption in aqueous mixtures containing chemical and physical solvents \cite{Cleeton2020,Bulow2021a}.
These studies motivate an electrolyte-EOS description of reactive absorption while retaining the solvent-specific chemistry, association choices, and parameter domains of each application.

% P:introduction-thermodynamic-gap
\Cref{tab:thermodynamic-model-comparison} summarizes the roles and input requirements of these descriptions.
The column question is how their equilibrium predictions propagate through interfacial transfer and the energy balance to capture and axial temperature.
The ePC-SAFT--eNRTL comparison addresses that question using two nonideal liquid descriptions; Henry's law supplies a simpler solubility reference.
Within ePC-SAFT, parameter sensitivity then identifies which specified thermodynamic inputs most influence the predicted column response.

\begin{table}[pos=htbp]
\centering\small
\caption{Thermodynamic descriptions relevant to reactive aqueous amine absorption, their uses, and their principal requirements. Model structure alone does not establish an accuracy ranking.}
\label{tab:thermodynamic-model-comparison}
\begin{tabularx}{\linewidth}{@{}>{\raggedright\arraybackslash}p{0.16\linewidth}>{\raggedright\arraybackslash}X>{\raggedright\arraybackslash}X@{}}
\toprule
Description & Thermodynamic role and practical use & Required inputs and limitations\\
\midrule
Henry-law reference & Molecular CO$_2$ solubility; a compact reference for the driving force & Solvent composition, temperature correlation, concentration basis, and a separate reaction system; applicability follows the correlation range\\
Kent--Eisenberg & Simplified reaction equilibrium; a compact correlation for represented solvent and loading ranges & Empirical equilibrium constants are solvent- and range-specific; electrolyte nonideality is treated approximately \cite{Shahid2019,Zhu2022}\\
eNRTL & Excess Gibbs energy and electrolyte activities; a flexible description of aqueous-amine VLE and speciation & Species, reaction standard states, and calibrated interaction parameters; extrapolation depends on the calibration domain \cite{Zhang2011}\\
CPA & Cubic equation of state with association; a common EOS description of physical and associating contributions & Physical, association, and binary parameters with the selected reaction description; reactive applications require additional chemical assumptions \cite{Chen2023}\\
ePC-SAFT & Residual Helmholtz energy with molecular and ionic contributions; a common basis for fugacity, density, and thermodynamic derivatives & Consistent segment, association, ionic, dielectric, and reaction-reference inputs; the EOS does not determine the reaction network or equilibrium constants \cite{Gross2001,Cameretti2005,Bulow2021a}\\
\bottomrule
\end{tabularx}
\end{table}

\subsubsection{Reactive transfer and numerical solution}
% P:introduction-08
Enhancement-factor models represent the interaction of reaction and diffusion through a correction to physical liquid-film transfer.
Gaspar and Fosb{\o}l developed a general enhancement-factor treatment for reversible absorption and desorption and applied it to the CO$_2$--MEA--water system \cite{Gaspar2015}.
Such approximations connect reaction kinetics to a compact transfer expression that can be evaluated throughout a column.
An equilibrium-manifold formulation instead considers the rapid-local-reaction limit and evaluates constrained mobility along a path of equilibrium states.
Comparing these descriptions at the same thermodynamic model and feeds tests the consequence of the transfer approximation for capture and the spatial distribution of absorption.
The comparison also separates the influence of species mobility from the physical film thickness.

% P:introduction-09
Published rate-based absorber models couple the transfer description to countercurrent material and energy balances \cite{Freguia2003,Zhang2009,Mores2012}.
Shooting, finite differences, and collocation provide different numerical representations of this boundary-value problem.
Their computational cost depends on scaling, initialization, resolution, and the thermodynamic work required at each local evaluation.
A comparison on the same physical equations and at matched accuracy connects numerical efficiency to the absorber calculation being performed.
This numerical comparison supports interpretation of the thermodynamic and film effects through a common resolution of capture and temperature.

% P:introduction-12
The literature establishes the coupled physical descriptions needed for MEA absorption, but agreement with property data alone does not determine the resulting packed-column response.
Comparisons that change thermodynamics, transport, numerical resolution, and operating cases together cannot assign a capture or temperature difference to one specified modeling choice.
The present study uses a common absorber formulation to compare thermodynamic descriptions and numerical methods, then examines the associated parameter sensitivity and liquid-film response.
The operating-condition study relates those controlled comparisons to local changes in feeds and properties.

\subsection{Scope and contributions}
% P:introduction-10
The central question is when additional thermodynamic and film detail materially changes or improves predicted capture and axial temperature, and what numerical effort is required to resolve that difference.
Three controlled studies address this question.
The thermodynamic study compares the nine-species ePC-SAFT and six-species eNRTL descriptions of equilibrium driving force and column response, with Henry's law as a reference.
Thermodynamic parameter sensitivity follows the model comparison and identifies the influence of four perturbation coordinates within ePC-SAFT.
The film study compares enhancement-factor and equilibrium-manifold transport with the thermodynamic model, feeds, and other constitutive laws held fixed.
The numerical study compares conservative simultaneous, shooting, finite-difference, and collocation formulations of the same coupled balances.
A Latin hypercube design then examines thermodynamic, transport, and operating inputs together to relate the controlled comparisons to the local operating response.
One-bed NCCC measurements provide the multi-case capture and axial-temperature evaluation, with Case 3C serving as the reference for local sensitivity \cite{Morgan2018a,Morgan2020}.

% P:introduction-13
The \emph{ePC-SAFT fugacity benchmark} compares liquid thermodynamic descriptions within a common absorber calculation.
Feed conditions, transport and hydraulic correlations, gas fugacity conventions, and numerical accuracy targets are held common within each comparison.
Each thermodynamic description retains its specified species, reaction constants, and reference states.
A comparison that changes those chemical inputs measures their combined effect with the activity or fugacity model; an isolated closure effect requires common species and reaction thermodynamics.
The common-property calculation holds density and caloric descriptions fixed, while the full-property calculation additionally includes the selected model's density and enthalpy response.
Predictive column performance means agreement with absorber observations excluded from parameter fitting; thermodynamic and transport inputs may be fitted independently to property measurements.

% P:introduction-11
The model substitutions define controlled comparisons conditional on the other model choices: interacting equilibrium, transport, and thermal effects do not have a unique additive attribution.
The relative performance of ePC-SAFT and the comparison models is assessed from capture and temperature together, without presupposing a model ranking.
Capture and axial liquid-temperature measurements provide the experimental comparisons, while material, energy, and charge conservation define the physical consistency of the calculation.
This structure connects the choice of absorber model to the separation and thermal predictions needed for post-combustion capture design.
% CHECKLIST-END: introduction
